
\section{ Motivation and Background}
We   argue here that  for ecosystems where verification is a primary bottleneck and data is scarce (like software engineering), the minimalist path of EZR.py is not a compromise, but a statistically and economically justified choice.


\subsection{Why not ``data-heavy''?} \label{dataheavy}
``Data-heavy'' models are often build in AI; e.g.
the large language models
built from terabytes of data.   But
EZR.py's data-lite approach reaches state-of-the-art results with
fewer than 100 labels (see below). This is
 useful since,  in SE: 
 \begin{itemize}
 \item
 {\bf  Labeling by experts} is sluggish and fatigue-prone \cite{easterby1980design}, taking hours per case for security or technical debt tasks.  
 \item Taking
 labels from {\bf historical logs} can be   unreliable; e..g   90\% of technical debt entries are false positives \cite{yu2020identifying}.  
 \item
{\bf Automated labeling oracles} can be wrong or slow to scale due to        approximate and   incorrect heuristics \cite{kamei2012large,DBLP:conf/wosp/ValovPGFC17}. 
 \end{itemize}
Aside: asking an LLM to label data does not resolve the labelling problem. Recent studies show human-model agreement is not a proxy for ground truth \cite{ahmed2024can}, meaning LLMs remain conditional collaborators rather than replacements.


\begin{table}[!t]
\caption{Examples from  SE. For $y=f(x)$,  acquiring $x$ inputs can be  
easier than finding the associated $y$ labels.}
\label{slow}
\begin{center}
{\footnotesize
\begin{tabular}{p{.44\columnwidth}|p{.44\columnwidth}}
\rowcolor{gray!20} Cheap feature Acquisition ($x$)  & Expensive label Determination ($y$)\\\hline
At scale, GitHub metadata (code size, dependency counts, etc.). & 
Estimate market valuation or total development man-hours.\\
\rowcolor{gray!10} Count system classes. & 
Secure permissions to audit actual maintenance effort.\\
List design option permutations. & 
Validate configurations via human stakeholder consensus.\\
\rowcolor{gray!10} Enumerate software parameters. & 
Execute full test suites for every unique build.\\
Identify hyper-parameters for data miners. & 
Perform exhaustive grid searches for local optima.\\
\rowcolor{gray!10} Generate fuzzing inputs. & 
Execute tests and audit outputs.\\
\end{tabular}}
\end{center}
\end{table}

\subsection{Why not ``code-heavy''?}
In this era of ``vibe coding,'' AI coding tools tempt us to build ``code-heavy'' applications hastily assembled from many packages . EZR.py instead uses extensive refactoring to minimize dependencies and total lines of code.

Such code minimization is good science. 
In his textbook on {\em Empirical AI}  \cite{cohen1995}, Cohen recommends baselining sophisticated methods against simpler straw-man implementations via {\em ablation studies}\footnote{Try removing or simplifying   parts of your system.}. EZR.py is the product of many  such ablations. We discarded multi-generational genetic algorithms by over-sampling in generation one \cite{nair2016accidental,chen18}. We found simple greedy search matches complex adaptive active learning sampling \cite{ganguly2025lowgodatalightse}. We replaced heavyweight packages like \texttto{argparse,pandas} with native Python (e.g.  our CSV and CLI readers are only 5 and 12 lines long, respectively).

Code minimization is good engineering. Vibe-coding works only if developers never look at the details—but that's exactly how you miss structural similarities across distant parts of a codebase. EZR.py (below) shows AI systems are full of such similarities. Spotting them lets us refactor and simplify, which cuts the downstream cost of review, maintenance, and porting. Developers who only vibe never get close enough to notice.

Another reason to minimize the code is to decrease  dependencies 
on third-party packages.   Modern software is   assembled from pre-made parts where 85--97\% of code is reused from external sources \cite{martinez2021}. While libraries like \texttt{pandas} are battle-tested, they introduce very large  ``attack surface'' for hackers
where a single poisoned update 
can silently compromise tens of thousands of other
software packages   without the 
primary developer's knowledge\cite{martinez2021}.
For this reason, EZR avoids packages like \texttt{pandas} in favor
of much smaller homer grown functions like \textsf{csv} (see line~\ref{csv} of the code).

\begin{table}[!b]
\caption{AI-Related Challenges in Software Engineering. }\label{problems}
\centering
\begin{tabular}{lp{6cm}}
\rowcolor{gray!20}\textbf{Issue} & \textbf{Evidence} \\\hline  
Cognitive debt & Generative AI acceleration can lead to teams
losing shared understanding of system design \cite{CognitiveDebt} \\
\rowcolor{gray!10}Hallucinations & 52\% of ChatGPT answers are incorrect (in part)\cite{S12} \\
Human error & Humans overlook AI's errors 39\% of the
time \cite{S12} \\
\rowcolor{gray!10}Code quality & 70\% of generated code is faulty: 29\% correct,
51\% partial, 20\% incorrect \cite{S15} \\
Expert caution & Developers exhibit skepticism, often accepting 
only a fraction of AI-suggested code \cite{S10} \\
\rowcolor{gray!10}Training bias & LLMs frequently regenerate simple, "stupid" bugs
present in their training data \cite{S13} \\
``AI Slop'' &   Individual productivity 
leads to technical debt\cite{S14} 
where costs are  externalize costs as  reviewers, maintainers, and the broader community
struggle to handle the ``AI slop''
generated from vibe coding~\cite{baltes2026anendlessstreamai}.
\end{tabular}

\end{table}\begin{table*}[!t] 
{\fontsize{9}{10}\selectfont % Optimal size between tiny and footnotesize
\centering
\setlength{\tabcolsep}{2.5pt} 
\renewcommand{\arraystretch}{1.1} 

\caption{Experimental Datasets: 124 benchmarks. Optimization objectives ($y$), features ($x$), and specific identifiers used in this study.}\label{mootdata}

\begin{tabularx}{\textwidth}{@{} r |  l l l X c | c | c | c @{}} 
\rowcolor{blue!15} 
\textbf{\#} & \textbf{Category} & \textbf{Focus} & \textbf{File Name(s)} & \textbf{Here, optimizers struggle with issues like...} & \textbf{\#y} & \textbf{\#x} & \textbf{Rows} & \textbf{Refs} \\ \midrule

% --- GROUP 1 ---
\multicolumn{9}{l}{\cellcolor{gray!30}\textbf{Software Systems \& Configuration \& Tuning}} \\
24 & Soft. Config & \textbf{PLE} & SS-[A-X] & Minimize footprint/memory vs. maximizing throughput & 2-3 & 3--88 & 197–86k & \cite{Amiraliminimaldata} \\ \rowcolor{gray!10}
12 & PromiseTune & \textbf{Perf.} & 7z, LLVM, BDBC & Minimize execution time and energy consumption & 1 & 6--35 & 864–166k & \cite{chen2026promisetune} \\ 
1 & Cloud compute& \textbf{Tuning} & Apache\_AllMeas & Balance server response vs. CPU/RAM load patterns & 1 & 9 & 192 & \cite{senthilkumar2024can} \\ \rowcolor{gray!10}
1 & Cloud  compute& \textbf{Tuning} & SQL\_AllMeas & Minimize latency/IO vs. maximizing throughput & 1 & 39 & 4,654 & \cite{chen2025accuracy} \\ 
1 & Cloud compute & \textbf{Tuning} & X264\_AllMeas & Optimize encoding parameters for PSNR/SSIM quality & 1 & 16 & 1,153 & \cite{Amiraliminimaldata} \\ \rowcolor{gray!10}
2 & Testing & \textbf{Testing} & test120, 600 & Maximize coverage while minimizing execution time & 1 & 9 & 5,161 & -- \\ \midrule

% --- GROUP 2 ---
\multicolumn{9}{l}{\cellcolor{gray!30}\textbf{Project Management \& Process Modeling}} \\
35 & Proj. Health & \textbf{Health} & Health-Closed & Optimize PR rates and minimize developer churn & 2-3 & 5 & 10,001 & \cite{lustosa2024learning} \\ \rowcolor{gray!10}
3 & Scrum & \textbf{Agile} & Scrum[1k-100k] & Maximize velocity within sprint constraints & 3 & 124 & 1k–100k & \cite{lustossa2024isneak} \\ 
8 & Feature Mod. & \textbf{Config} & FFM, FM-* & Optimize Clause/Constraint ratio in large spaces & 3 & 128--1k & 10,001 & \cite{Amiraliminimaldata} \\ \rowcolor{gray!10}
1 & nasa93dem & \textbf{Cost} & nasa93dem & Minimize effort (person-months) vs. quality & 3 & 26 & 93 & \cite{lustosa2024learning} \\ 
1 & COC1000 & \textbf{Risk} & coc1000 & Minimize risks/schedule slip vs. analyst expertise & 5 & 20 & 1,001 & \cite{chen2018beyond} \\ \rowcolor{gray!10}
4 & POM3 & \textbf{Process} & pom3[a-d] & Balance project idle rates vs. completion costs & 3 & 9 & 501–20k & \cite{lustossa2024isneak} \\ 
4 & XOMO & \textbf{Defects} & xomo\_[flt,grd] & Minimize defect density vs. total project effort & 4 & 27 & 10,001 & \cite{chen2018beyond} \\ \midrule

% --- GROUP 3 ---
\multicolumn{9}{l}{\cellcolor{gray!30}\textbf{Interdisciplinary \& Behavioral Benchmarks}} \\ \rowcolor{gray!10}
4 & Behavioral & \textbf{Behavior} & all\_players, etc & Maximize user retention and predict cohort churn & 1-3 & 26--55 & 82–17k & \cite{nyagami_fc25_kaggle_2025} \\ 
4 & Financial & \textbf{Finance} & BankChurners & Minimize credit risk vs. optimal pricing models & 2-5 & 19--77 & 1k–20k & \cite{blastchar_telco_customer_churn_2025} \\ \rowcolor{gray!10}
3 & Health & \textbf{Health} & COVID19, Life & Maximize accuracy/life vs. readmission likelihood & 1-3 & 20--64 & 2k–25k & \cite{kumarajarshi_life_expectancy_who_2025} \\ 
2 & RL & \textbf{Control} & A2C\_[Acr,Crt] & Maximize reward signals vs. steps to convergence & 3-4 & 9--11 & 224–318 & -- \\ \rowcolor{gray!10}
5 & Sales & \textbf{Market} & Marketing, socks & Maximize ROI vs. minimizing forecasting error & 1-8 & 20--31 & 247–2k & \cite{jackdaoud_marketing_data_2022} \\ 
3 & Misc. & \textbf{General} & auto93, Car & Optimize multivariate trade-offs (e.g. MPG vs. HP) & 2-5 & 5--38 & 205–1k & \cite{Amiraliminimaldata} \\ \midrule

\rowcolor{blue!15}
\textbf{124} & \multicolumn{8}{l}{\textbf{Total}}  
\end{tabularx}
\end{table*}


``Backpacking'' is the process of cutting back the code to only what is most useful\cite{hipp2021}.
Minimizing dependencies, with backpacking  helps software remain functional, secure, and independent. For example,
backpacking meant that SQLite author Richard Hipp avoided external vendors by building his own editor, servers, and storage engine. By shunning Berkeley DB
(which Oracle later acquired and then demanded license fees)
he ensured SQLite remained free. Consequently, SQLite now ships on all Android/iOS devices, major browsers (Chrome, Firefox, Safari), most Linux distributions, and even Mars rovers.  As Hipp puts it:
{\em ``If you want to be free, that means doing things yourself''}\cite{hipp2021}.


\subsection{Why not ``model-heavy''?}
EZR.py runs in milliseconds to generate. An alternate ``model-heavy'' approach, advocated in Sutton’s famous {\em Bitter Lesson} essay\footnote{\url{http://www.incompleteideas.net/IncIdeas/BitterLesson.html}}, is to always
 leverage computation at scale since, as Sutton  says  ``general methods that leverage computation are
ultimately the most effective''.

For several reasons, we take issue with  Sutton's   approach.
{\bf First}, mega-scale compute may be
unavailable or inappropriate. Many edge devices operate under power constraints where algorithmic efficiency is a necessity, not a distraction. Also, when there is very little domain data (as seen in SE), mega-computation is superfluous.


{\bf Second}, It may be unwise to assume  compute budgets can grow without limit. Analysts warn of an \$800 billion gap between AI revenue and data-center costs through 2030. Training data is running out \cite{villalobosposition}, and training on the resulting synthetic substitutes risks "model collapse" \cite{shumailov2024}. Clearly, the era of  "more is better"   is hitting a wall.

 


{\bf Third}, Sutton’s arguments for mega-computation used
examples with clear objective functions like vision or gaming. In SE,
initial success  with LLMs often has subsequent hidden costs like security flaws
\cite{S15} and cognitive debt \cite{CognitiveDebt} (see
Table~\ref{problems}). While next-generation models might resolve these
issues, scaling AI accelerates the production of ``AI slop''
\cite{baltes2026anendlessstreamai}, shifting effort to a grueling
verification bottleneck \cite{S10}. This creates a paradox: as AI
output grows, so does the human burden of checking it.


{\bf Fourth}, safety-critical domains require auditable, repairable models. 
Krikorian~\cite{krikorian2026tesla} warns  that   humans
tend  to switch off their critical
faculties when a system, like AI, usually performs satisfactorily. This means,  when AI inevitably falters, humans either do not notice or cannot react in a timely manner.
As Rudin \cite{rudin2019stop} argues, we should use interpretable models (like those generated by EZR.py) instead of struggling (or overlooking) with the need  to explain black boxes.


{\bf Fifth}, the complexities and nuances explored by a data-heavy
 In such domains, below an epsilon (e.g. Cohen's \mbox{$d < 0.35\sigma$}), further “improvements” are actually spurious. 


